\chapter{Results: The Arena}\label{ch:arena}

\begin{center}
\includegraphics[width=0.56\linewidth]{figures/arena_board.png}
\captionof{figure}{An arena game at turn 120: six copies of the arena network on a
$19\times19$ board, all six still alive.}
\label{fig:arena-board}
\end{center}

The arena is a free-for-all: six snakes on a $19\times19$ board, with at least six food
on the board at all times, so a crowded board does not simply starve. The arena network
reads the $17\times17$ window of Section~\ref{sec:window}, so the same weights also play
the browser's arena with anywhere from 2 to 8 snakes.

\section{Measuring a six-player game}

\begin{keybox}[One against five]
A six-player game cannot be scored as a single head-to-head. Instead the network takes
one seat and the other \emph{five} are all the same scripted bot. If the network were
exactly as good as the bot, it would win $1/6\approx16.7\%$ of games. Every arena win rate
should be read against that fair share, not against 50\%.
\end{keybox}

\section{Learning curve}

\begin{center}
\begin{tikzpicture}
\begin{axis}[reportaxis, xmin=0, xmax=16.5, ymin=0, ymax=100,
  xlabel={self-play game turns (millions)}, ylabel={win rate as 1 of 6 snakes (\%)},
  legend style={at={(0.5,-0.24)},anchor=north,legend columns=2,/tikz/every even column/.append style={column sep=1.2em}}]
\addplot[color=cmured] table[x=steps, y=flood, col sep=comma] {data/arena_eval.csv};
\addlegendentry{vs $5\times$ flood-fill}
\addplot[color=inkblue] table[x=steps, y=greedy, col sep=comma] {data/arena_eval.csv};
\addlegendentry{vs $5\times$ greedy}
\addplot[color=rulegrey, dashed, line width=0.9pt, forget plot] coordinates {(0,16.67) (16.5,16.67)};
\node[font=\scriptsize\itshape, color=cmudark, anchor=north east] at (axis cs:16.2,15.5) {fair share, $1/6$};
\addplot[color=amber, densely dotted, line width=1pt, forget plot] coordinates {(8.847,0) (8.847,100)};
\node[font=\scriptsize\itshape, color=amber!80!black, anchor=north west, align=left] at (axis cs:9.0,30) {highest\\60-game score};
\end{axis}
\end{tikzpicture}
\captionof{figure}{Arena win rate during training, evaluated every 30 iterations over 60
games. The dotted line marks the checkpoint that scored highest on these small
evaluations; it was not the one finally used (see below).}
\label{fig:arena-curve}
\end{center}

Against five greedy bots the network is ahead of its fair share from the very first
checkpoint, at 77\% after half a million turns. Greedy bots in a crowd mostly kill each
other, and a snake that merely avoids them survives. Against five flood-fill bots it
starts below fair share, at 8\%, and passes 40\% within 4M turns.

\begin{center}
\begin{tikzpicture}
\begin{axis}[reportaxis, height=4.8cm, xmin=0, xmax=16.5, ymin=0,
  xlabel={self-play game turns (millions)}, ylabel={mean game length}]
\addplot[color=cmured] table[x=steps, y=turns, col sep=comma] {data/arena_train.csv};
\end{axis}
\end{tikzpicture}
\captionof{figure}{Mean length of six-snake self-play games, growing from about 50 turns
to over 1{,}000 as the snakes learn to share the board without colliding.}
\label{fig:arena-length}
\end{center}

\section{Choosing the checkpoint}

\begin{warnbox}[The ``best'' checkpoint was not the best]
Checkpoints were scored during training on only 60 games each. Picking the maximum of 33
noisy scores favours a checkpoint that got lucky, a selection bias. The highest scorer
was iteration 540, at 8.8M turns, yet later checkpoints looked at least as strong. Two
candidates were therefore compared on a separate seed, 777, never used for the final
numbers:
\begin{center}
\begin{tabular}{@{}lcc@{}}
\toprule
\textbf{200 games, seed 777} & \textbf{vs $5\times$ flood-fill} & \textbf{vs $5\times$ greedy}\\
\midrule
iteration 540 (highest 60-game score) & 35.0\% & 81.0\%\\
iteration 971 (end of training) & \textbf{47.5\%} & \textbf{92.5\%}\\
\bottomrule
\end{tabular}
\end{center}
The final checkpoint was chosen. The reported results below come from a third seed,
12345, so neither the training-time scores nor this comparison inflates them.
\end{warnbox}

\section{Final evaluation}

\begin{center}
\begin{tabular}{@{}lcccccc@{}}
\toprule
\textbf{Opponents} & \textbf{W} & \textbf{D} & \textbf{L} & \textbf{Win rate (95\% CI)} &
  \textbf{Game length} & \textbf{Final length}\\
\textbf{($\times5$)} & & & & & \footnotesize turns & \footnotesize net\ /\ bots\\
\midrule
flood\_fill & 134 & 2 & 164 & \textbf{44.7\%} \ (39.1--50.3) & 385.2 & 43.0 / 25.1\\
greedy & 264 & 2 & 34 & \textbf{88.0\%} \ (83.8--91.2) & 163.8 & 24.7 / 14.3\\
safe\_random & 296 & 0 & 4 & \textbf{98.7\%} \ (96.6--99.5) & 205.7 & 32.4 / \hphantom{1}5.3\\
\bottomrule
\end{tabular}
\end{center}
\begin{center}
\footnotesize\color{cmudark} 300 games per row; fair share $16.7\%$.
\end{center}

\begin{keybox}[Nearly three times its fair share]
Against a table of five flood-fill bots, each built to protect its own space, the network
wins 44.7\% of six-player games, $2.7\times$ the $16.7\%$ an equal player would get. It
also grows much longer than its opponents, 43 segments against 25. In a crowded game,
length is both a weapon (it wins head-to-heads) and a wall that takes space from
everyone else, and the network uses it as both.
\end{keybox}

\section{The window generalises}

Because the network never sees the whole board, nothing in it depends on the board size
or the number of snakes. In the browser it plays two to eight snakes on $19\times19$
without retraining. No numbers are reported for configurations other than the
six-snake training setup, so how well this generalisation holds up is unmeasured.
